% Experiments --- setup, datasets, metrics, harness. Numbers live in
% Results; this section fixes the protocol. The LLM-benchmark table is a
% honest placeholder pending US-069 (H100 eval), per the data-real rule.
\section{Configuración Experimental}
\label{sec:experiments}

\subsection{Conjuntos de datos}
\paragraph{PASTIS-R}~\cite{garnot2021pastis} es el benchmark principal:
2{,}433 parches Sentinel-2 multitemporales ($128\times128$ px), 18 clases
de cultivo, Francia metropolitana, con cinco folds espacialmente
disjuntos. Salvo que se indique lo contrario, los modelos de
segmentación se evalúan en un fold de prueba retenido y los combinadores
de ensamble sobre predicciones out-of-fold.

\paragraph{Sen4AgriNet (Cataluña)}~\cite{sykas2022sen4agrinet}
(CC-BY-SA-4.0) se usa para el estudio de transferencia
Francia$\rightarrow$Cataluña. Un subconjunto de Cataluña se versiona con
DVC (\texttt{data/sen4agrinet.dvc}).

\paragraph{EuroCropsML}~\cite{reuss2025eurocropsml} (CC-BY-SA-4.0)
sustenta la curva few-shot de $k$-shots; el subconjunto crudo se versiona
con DVC (\texttt{data/transfer/eurocropsml.dvc}).

\paragraph{AgroMind}~\cite{ruan2025agromind} se usa solo para evaluar la
capa conversacional. No tiene partición de entrenamiento (alrededor de
28{,}482 pares de QA); el ajuste fino sobre él induciría fuga, por lo que
es estrictamente un conjunto de evaluación.

\subsection{Métricas}
La segmentación se reporta con la intersección sobre la unión media
(mIoU) y el F1 promediado en macro a nivel de parcela; también
reportamos la exactitud de píxel donde sea relevante. El promedio en
macro es el objetivo principal porque penaliza por igual los errores en
cultivos raros y comunes, lo que coincide con el objetivo de negocio.
Los combinadores de ensamble se comparan por F1 macro y exactitud
out-of-fold.

\subsection{Arnés de evaluación y reproducibilidad}
Todas las corridas de segmentación se puntúan con un arnés de evaluación
compartido (\texttt{ml/eval/}) que impone folds espacialmente disjuntos
y aborta ante un solapamiento geográfico entre los conjuntos de
entrenamiento y evaluación. Un módulo de remapeo de etiquetas
(\texttt{ml/eval/class\_remap.py}) proyecta las predicciones sobre el
espacio macro de etiquetas armonizado para la comparación entre
regiones. Cada cifra reportada en la Sección~\ref{sec:results} está
anclada a un artefacto versionado (JSON/CSV/registro) en el repositorio,
registrado en comentarios \texttt{\% src:} en el código fuente; los
conjuntos de datos y los pesos se rastrean con DVC y las corridas con
MLflow.

\subsection{Protocolo de benchmark de la capa conversacional}
\label{sec:experiments-llm}
La capa conversacional se evalúa sobre AgroMind con un protocolo de
LLM-como-juez que compara el razonador congelado Gemini~2.5~Pro frente al
razonador multimodal Qwen on-premises servido vía vLLM. La tabla
cuantitativa del benchmark se produce mediante una corrida de evaluación
separada en la GPU H100 (US-069) y se deja aquí intencionalmente como un
marcador de posición estructurado; las cifras no se fabrican.
La Tabla~\ref{tab:llm-bench} fija la estructura de columnas.

\begin{table}[htbp]
\caption{Benchmark de la capa conversacional (solo estructura; cifras
pendientes de la evaluación en H100, US-069).}
\label{tab:llm-bench}
\centering
\begin{tabular}{lccc}
\toprule
\textbf{Razonador (congelado)} & \textbf{Exactitud} & \textbf{Anclaje}
  & \textbf{LLM-juez} \\
\midrule
Gemini~2.5~Pro (nube)       & --- & --- & --- \\
Qwen3.5-35B-A3B (on-prem)   & --- & --- & --- \\
\bottomrule
\end{tabular}
% src: pending US-069 (H100 eval) -- do not fill with synthetic numbers.
\end{table}
